\chapter*{Notation \& Conventions}
\addcontentsline{toc}{chapter}{Notation \& Conventions}

\section{Typographic Conventions}

This manual uses consistent typographic formatting to distinguish between
different kinds of information.  Table~\ref{tab:typographic-conventions}
provides a summary.

\begin{table}[htbp]
    \centering
    \caption{Typographic conventions used in this manual}
    \label{tab:typographic-conventions}
    \begin{tabular}{lll}
        \toprule
        \textbf{Appearance}      & \textbf{Meaning}                    & \textbf{Example} \\
        \midrule
        \texttt{Monospace}       & \LaTeX{} commands, code, filenames & \texttt{\textbackslash documentclass} \\
        \textsf{Sans-serif}      & Key-value options, package options  & \textsf{language=english} \\
        \textbf{Bold}            & Key terms (first occurrence)        & \textbf{doctype profile} \\
        \fbox{\small UI element} & Menu paths, interface elements      & \fbox{\small File\;\textgreater\;Compile} \\
        \midrule
        \emph{Italics}           & Emphasis, introduced terminology    & \emph{preamble} \\
        \textsc{Small Caps}      & Proper names of standards           & \textsc{Unicode} \\
        \bottomrule
    \end{tabular}
\end{table}

When a command is shown with its mandatory argument, curly braces are included
(e.g.\ \texttt{\textbackslash chapter\{Title\}}).  Optional arguments appear
in square brackets (e.g.\ \texttt{\textbackslash section[short]\{Long Title\}}).

\section{Package References}

Throughout this manual, \LaTeX{} packages are referenced using the
\texttt{\textbackslash ctanpackage\{\}} command, which produces a hyperlink to
the package's page on the Comprehensive \TeX{} Archive Network (CTAN).  For
example, the \ctanpackage{minted} package is referenced as shown.  Hovering
over or clicking the link navigates to \url{https://ctan.org/pkg/minted}.

This convention allows readers to quickly locate the authoritative
documentation for any package mentioned in the text.

\section{Cross-References}

The manual uses the standard \LaTeX{} cross-referencing system.  Chapters are
numbered consecutively within each part.  Equations, figures, tables, and
listings are numbered per chapter using the scheme
\texttt{<chapter>.<counter>}.  For example, Equation~(3.7) refers to the
seventh equation in Chapter~3.

Cross-references within the manual use the commands
\texttt{\textbackslash ref\{\}}, \texttt{\textbackslash eqref\{\}},
\texttt{\textbackslash cref\{\}}, and \texttt{\textbackslash Cref\{\}} from
the \ctanpackage{cleveref} package.  Readers should note that
\texttt{\textbackslash cref} and \texttt{\textbackslash Cref} automatically
insert the appropriate label prefix (``Figure'', ``Table'', ``Chapter'', etc.).

\section{Code Examples}

All code examples presented in this manual are intended to be directly
compilable when placed in a suitable document preamble.  Code listings are
rendered using the \ctanpackage{minted} package, which provides syntax
highlighting for a wide range of programming and markup languages.

Where an example is shortened for brevity, the ellipsis marker
\texttt{\ldots} is used within the listing and the omission is explicitly
described in the surrounding text.  Complete source files for the examples are
available in the \texttt{examples/} directory of the OmniLaTeX repository.

\section{Glossary Symbols}

OmniLaTeX integrates the \ctanpackage{glossaries-extra} package to provide
consistent symbol and abbreviation management.  The following shorthand
commands are used throughout this manual:

\begin{description}
    \item[\texttt{\textbackslash sym\{\}}] References a mathematical symbol
        entry from the glossary.  For instance,
        \texttt{\textbackslash sym\{velocity\}} produces the symbol and its
        associated description as defined in the glossary file.
    \item[\texttt{\textbackslash abb\{\}}] References an abbreviation entry.
        First use expands the abbreviation; subsequent uses display the short
        form.  For example, \texttt{\textbackslash abb\{LaTeX\}} expands to
        ``\LaTeX{}'' on first use and ``\LaTeX{}'' thereafter.
    \item[\texttt{\textbackslash num\{\}}] Formats a number according to the
        locale settings of the document.  This is particularly useful in
        multilingual documents where number formatting conventions differ.
\end{description}

The glossary definitions themselves are maintained in dedicated
\texttt{.tex} files under the \texttt{glossaries/} directory and are loaded
via \texttt{\textbackslash input\{\}} in the document preamble.  Detailed
instructions for creating and managing glossary entries can be found in
Chapter~\ref{ch:glossaries}.
